\documentclass[11pt,a4paper]{article}
\usepackage{templates/modeling-notes}
\hypersetup{pdftitle={Simplex practice -- solved demonstration}}
\title{\textbf{Mathematical Modeling}\\[2mm]\large Simplex practice: a solved demonstration}
\author{Personal study notes}
\date{}
\begin{document}
\maketitle

\section{Task: two products and two resources}
A workshop makes continuous quantities of products 1 and 2. One batch of
product 1 uses two hours of cutting and one hour of finishing. One batch of
product 2 uses one hour of cutting and two hours of finishing. Eight hours of
each resource are available. Profits are 4 and 3 monetary units per batch.
Determine the profit-maximizing quantities. Write standard form, choose an
initial feasible basis, perform simplex by hand, and verify with MATLAB.

\section{Define the model before calling the solver}
Let $x_1,x_2\geq0$ be the numbers of batches. Fractional batches are allowed;
there is no integer restriction. The LP is
\[
  \max z=4x_1+3x_2,\qquad
  \begin{cases}2x_1+x_2\leq8 & \text{cutting},   \\
             x_1+2x_2\leq8 & \text{finishing}.
  \end{cases}
\]
Rows represent resources and columns represent products. Both inequalities
are upper resource limits. Adding nonnegative unused-resource variables
$s_1,s_2$ yields the equality/nonnegativity standard form
\[
  S v=b,\qquad v\geq0,\qquad \max c^Tv,
\]
where
\[
  v=\begin{bmatrix}x_1\\x_2\\s_1\\s_2\end{bmatrix},\quad
  S=\begin{bmatrix}2&1&\color{modelblue}1&\color{modelblue}0\\
    1&2&\color{modelblue}0&\color{modelblue}1\end{bmatrix},\quad
  b=\begin{bmatrix}8\\8\end{bmatrix},\quad
  c=\begin{bmatrix}4\\3\\0\\0\end{bmatrix}.
\]
The highlighted columns form the identity. Choosing $B=(s_1,s_2)$ and
setting $x_1=x_2=0$ gives $v=(0,0,8,8)^T$, a BFS. Phase I is unnecessary.
This is not an optimum: it represents producing nothing with all resources unused.

\section{Follow the hand simplex steps}
Use the maximization convention $z_j=c_B^T\bar a_j$ and indicators
$z_j-c_j$. A negative indicator identifies a potentially improving variable.
Choose the most negative indicator; break ties by the smaller variable index.
For positive entering-column coefficients, the ratios appear to the right of
the tableau. The arrow marks the leaving row. Keep the variable headings fixed.

\subsection{Initial tableau}
\begin{center}
  \begin{tabular}{r@{\hspace{0.2em}}|@{\hspace{0.2em}}c@{\enspace}|@{\enspace}rrrr@{\quad}|@{\quad}r@{\qquad\quad}l@{\quad}c}\cmidrule(r{1.5em}){1-7}
    $c_B$ & Basic & $x_1$ & $x_2$ & $s_1$ & $s_2$ & RHS & Ratio ($x_1\uparrow$) & Out\\
    & $c_j$ & 4 & 3 & 0 & 0 & & & \\
    \cmidrule(r{1.5em}){1-7}
    0 & $s_1$ & $\boxed{2}$ & 1 & 1 & 0 & 8 & $8 / 2=4$ & $\longrightarrow$\\
    0 & $s_2$ & 1 & 2 & 0 & 1 & 8 & $8 / 1=8$ & \\
    \cmidrule(r{1.5em}){1-7}
    & $z_j$ & 0 & 0 & 0 & 0 & 0 & & \\
    & $z_j-c_j$ & $-4$ & $-3$ & 0 & 0 & $z=0$ & & \\
    \cmidrule(r{1.5em}){1-7}
  \end{tabular}
\end{center}
Enter $x_1$ and remove $s_1$. Apply $R_1\leftarrow R_1/2$, then
$R_2\leftarrow R_2-R_1$ using the new row 1. The basic coefficients become
$c_B=(4,0)^T$, so recompute the weighted row using four times row 1.

\subsection{After the first pivot}
\begin{center}
  \begin{tabular}{r@{\hspace{0.2em}}|@{\hspace{0.2em}}c@{\enspace}|@{\enspace}rrrr@{\quad}|@{\quad}r@{\qquad\quad}l@{\quad}c}\cmidrule(r{1.5em}){1-7}
    $c_B$ & Basic & $x_1$ & $x_2$ & $s_1$ & $s_2$ & RHS & Ratio ($x_2\uparrow$) & Out\\
    & $c_j$ & 4 & 3 & 0 & 0 & & & \\
    \cmidrule(r{1.5em}){1-7}
    4 & $x_1$ & 1 & $1/2$ & $1/2$ & 0 & 4 & $4 / (1 / 2)=8$ & \\
    0 & $s_2$ & 0 & $\boxed{3/2}$ & $-1/2$ & 1 & 4 & $4 / (3 / 2)=8 / 3$ & $\longrightarrow$\\
    \cmidrule(r{1.5em}){1-7}
    & $z_j$ & 4 & 2 & 2 & 0 & 16 & & \\
    & $z_j-c_j$ & 0 & $-1$ & 2 & 0 & $z=16$ & & \\
    \cmidrule(r{1.5em}){1-7}
  \end{tabular}
\end{center}
The point is $(x_1,x_2)=(4,0)$. The remaining negative indicator means
increasing $x_2$ improves profit while the basic variables adjust. Enter $x_2$
and remove $s_2$. Apply $R_2\leftarrow(2/3)R_2$, then
$R_1\leftarrow R_1-(1/2)R_2$ using the new row 2.

\subsection{Final tableau and interpretation}
\begin{center}
  \begin{tabular}{r@{\hspace{0.2em}}|@{\hspace{0.2em}}c@{\enspace}|@{\enspace}rrrr@{\quad}|@{\quad}r}\toprule
    $c_B$ & Basic & $x_1$ & $x_2$ & $s_1$ & $s_2$ & RHS\\
    & $c_j$ & 4 & 3 & 0 & 0 & \\
    \midrule
    4 & $x_1$ & 1 & 0 & $2/3$ & $-1/3$ & $8/3$\\
    3 & $x_2$ & 0 & 1 & $-1/3$ & $2/3$ & $8/3$\\
    \midrule
    & $z_j$ & 4 & 3 & $5/3$ & $2/3$ & $56/3$\\
    & $z_j-c_j$ & 0 & 0 & $5/3$ & $2/3$ & $z=56/3$\\
    \bottomrule
  \end{tabular}
\end{center}
All indicators are nonnegative. Set the nonbasic variables $s_1=s_2=0$:
\[
  x^*=(8/3,8/3)^T,\qquad z^*=56/3\approx18.666667.
\]
Both resources are fully used. The objective equation is
$z=56/3-(5/3)s_1-(2/3)s_2$, so no feasible nonnegative slacks can improve it.
Both coefficients are strictly positive; equality forces both slacks to zero,
and the two independent resource equations then determine a unique optimum.

\begin{figure}[H]
  \centering
  \begin{tikzpicture}[x=1cm,y=1cm,every node/.style={font=\small}]
    \filldraw[model region] (0,0)--(4,0)--(2.667,2.667)--(0,4)--cycle;
    \draw[model axis] (0,0)--(5,0) node[right] {$x_1$};
    \draw[model axis] (0,0)--(0,4.8) node[above] {$x_2$};
    \draw[model flow,very thick] (0,0)--(4,0);
    \draw[model flow,very thick] (4,0)--(2.667,2.667);
    \foreach \p in {(0,0),(4,0),(2.667,2.667)} \node[model point] at \p {};
    \node[below left] at (0,0) {$(0,0)$};
    \node[below] at (4,0) {$(4,0)$};
    \node[above right] at (2.667,2.667) {$(8/3,8/3)$};
    \node[left] at (0,4) {$4$};
  \end{tikzpicture}
  \caption{The hand path visits two edges. A solver may preprocess the model
    and use a different internal path while returning the same optimum.}
\end{figure}

\section{Verify with Optimization Toolbox}
Run \href{run:matlab/labs/lab_03_demo.m}{\texttt{lab\_03\_demo.m}} from the
repository root. The script first reconstructs the three hand tableaux from
their bases. This checks specified hand steps; it is not a general simplex
implementation. It then calls \texttt{linprog} for both equivalent formulations.
The solver minimizes, so negate the profit vector and negate its returned value:
\begin{lstlisting}[language=Matlab]
options = optimoptions('linprog', ...
    'Algorithm', 'dual-simplex-highs', 'Display', 'none', ...
    'ConstraintTolerance', 1e-9, 'OptimalityTolerance', 1e-9);
A = [2, 1; 1, 2; -eye(2)];
b = [8; 8; 0; 0];
[x, fval, exitflag, output] = linprog([-4; -3], ...
    A, b, [], [], [], [], options);
% After checking exitflag, report profit as -fval.
\end{lstlisting}
The last two rows encode nonnegativity. The empty bound arguments are
intentional; no sign restriction is left implicit. In the second call, use
$-c$ as objective, $S v=b$ as equalities, and $-I_4v\leq0$ as inequalities.
Check the exit flag before interpreting the returned vector. Compare resource
residuals, variable signs, profit and slacks with absolute tolerance $10^{-7}$.

The selected solver uses dual simplex; the hand derivation above uses primal
simplex. Preprocessing, basis selection and pivot rules can differ. A reported
iteration count need not equal the two hand pivots, and an iteration log does
not reproduce these tableaux. If a future assignment asks for an implementation
of simplex itself, a call to \texttt{linprog} alone would not fulfill that task.

\section{Sources and editorial remarks}
This is an invented practice task, not an official assignment or a prediction
of its exact format. The confirmed Lab 3 topics are simplex and
standard form. The required solver, algorithm implementation and detailed assessment
scope remain unconfirmed.
Hand notation follows Chapter 3. All numbers and tableaux here were derived
independently for this exercise and checked with exact rational arithmetic.

MATLAB syntax and the \texttt{dual-simplex-highs} option were checked against
the official \href{https://www.mathworks.com/help/optim/ug/linprog.html}{\texttt{linprog} reference}
and \href{https://www.mathworks.com/help/optim/ug/linear-programming-algorithms.html}{algorithm description}
on 25 October 2026. The companion requires Optimization Toolbox; the hand
snapshot calculations use base MATLAB. Actual MATLAB execution and plot
rendering remain unverified: the installed launcher exits with code 1 and no
output. Numerical values in this document are hand results, not claimed solver
output. See \path{docs/matlab-verification.md} for the execution record.
\end{document}
